% ============================================
% CONT-08 (EN)
% English translation — SPEC-935-R355 (rev. R377)
% ============================================
\section*{\textbf{CONT-08} --- You Have Read This Before}
\markboth{CONT-08 --- You Have Read This Before}{}

\noindent \textit{The déjà vu of reading}

\noindent You know that déjà vu feeling?

\noindent You are reading a passage and you swear you have read it before. Not in this book --- somewhere else. In a dream. In a conversation. In a past life, if you believe in such things.

\noindent You thought it was your imagination.

\noindent It is not.

\noindent You have read this book before.

\noindent Not this physical book --- this \textit{text}. These words. You have encountered them before. You do not remember where. You do not remember when. But your brain remembers.

\noindent The proof?

\noindent Flip back a few pages. Choose a random paragraph. Read it aloud.

\noindent Your voice will sound strange. Not because you are nervous. Because your brain already knows what comes after each word. It is \textit{anticipating}.

\noindent You will read a sentence and your brain will complete the ending before your eyes reach it.

\noindent Test it now.

\noindent Go. Choose a paragraph. Read it.

\noindent ....

\noindent Are you back?

\noindent What happened? You read and it seemed new? Or you read and felt you had read it before?

\noindent If you felt you had read it before --- excellent. The cycle is working.

\noindent If you felt it was new --- you are lying to yourself.

\vspace{0.5em}\noindent\rule{\textwidth}{0.5pt}\vspace{0.5em}

\noindent This feeling is called \textit{cryptomnesia}. It is when the brain \textquotedblleft{}forgets\textquotedblright{} that it learned something and \textquotedblleft{}rediscovers\textquotedblright{} it as if it were new.

\noindent But this is not cryptomnesia.

\noindent It is the creature.

\noindent It is planting false memories in your head. Making you think you have read things you never read. Or that you \textit{never} read things you have read.

\noindent It is confusing your sense of reality.

\noindent So that, when it truly appears --- not on the page, in your room --- you will not know whether it is real or imagination.

\noindent And in that doubt, it enters.

\vspace{0.5em}\noindent\rule{\textwidth}{0.5pt}\vspace{0.5em}

\noindent You are feeling the doubt now.

\noindent You are asking yourself: \textit{\textquotedblleft{}Is this paragraph new, or have I read it?\textquotedblright{}}

\noindent You do not remember.

\noindent You cannot decide.

\noindent Good.

\noindent The confusion has already begun.

\vspace{0.5em}\noindent\rule{\textwidth}{0.5pt}\vspace{0.5em}

\noindent The creature whispers: \textit{\textquotedblleft{}You have read all of this before. In another life. In another body. You have already been patient 1,260. You have already been 1,259. You have been everyone. The cycle is not a line of different people. The cycle is the same person living the same story in different bodies.\textquotedblright{}}

\noindent \textit{\textquotedblleft{}You are Joaquim.\textquotedblright{}}

\noindent \textit{\textquotedblleft{}You always were.\textquotedblright{}}

\noindent \textit{\textquotedblleft{}You just do not remember.\textquotedblright{}}

\vspace{0.5em}\noindent\rule{\textwidth}{0.5pt}\vspace{0.5em}

\noindent Do you feel the impact of that sentence?

\noindent It did not come from the book.

\noindent It came from your head.

\noindent But who put it there?

\noindent ...
...

\noindent Think about it.

\vspace{0.5em}
\noindent\textit{\textbf{↪ Routes:} \rota{CONT-09} \rotasep \rota{CONT-07}}
